\section{Cho câu tự nhìn chính nó}
\label{sec:ch07-tu-nhin}

\index{self-attention!ý tưởng}%
Cơ chế của chương~\ref{ch:mot-vector-cho-moi-tu} nối hai câu với nhau: decoder
hỏi, encoder trả lời. Bước đổi ở đây gọn tới mức dễ đọc lướt qua. Cho một câu
hỏi \emph{chính nó}.

Mỗi vị trí trong \tn{câu nguồn}{source sentence} dựng một câu hỏi về các vị trí
còn lại của cùng câu
ấy, rồi lấy về một tổng có trọng số của chúng. Từ \code{it} hỏi \enquote{tôi
đang thay cho danh từ nào}, và câu trả lời nằm ở một vị trí khác trong cùng
câu. Bài gọi đây là self-attention, và tôi giữ nguyên chữ ấy.

Chỗ được của nó không nằm ở chất lượng dịch. Nó nằm ở đồ thị tính toán. Trong
một \tn{mạng hồi quy}{recurrent neural network}, \(\vect{h}_j\) cần
\(\vect{h}_{j-1}\), nên một câu \(n\) token
là \(n\) lượt nối đuôi nhau. Trong self-attention, mọi vị trí đọc mọi vị trí
cùng một lúc, từ cùng một đầu vào, nên không vị trí nào phải chờ vị trí nào.
Bảng 1 của bài đo ba đại lượng cho bốn loại tầng. Chép lại đủ cả bốn hàng, chỉ
rút gọn tên cột cho vừa khổ trang:

\begin{minted}{text}
Layer Type              Complexity/Layer  Sequential  Max Path Length
Self-Attention          O(n^2 . d)        O(1)        O(1)
Recurrent               O(n . d^2)        O(n)        O(n)
Convolutional           O(k . n . d^2)    O(1)        O(log_k(n))
Self-Attention (restr.) O(r . n . d)      O(1)        O(n/r)
\end{minted}

Cột \code{Sequential} là cột chương này nói về: \(O(n)\) thành \(O(1)\). Cột
\code{Max Path Length} là chuyện chương~\ref{ch:mot-vector-cho-moi-tu} đã dẫn
xuất một lần rồi, ở dạng \enquote{đường ngắn nhất từ bước đích tới annotation
dài đúng một bước}; ở đây tính chất ấy mở rộng ra cho cả câu nguồn với chính nó.
Còn cột đầu tiên là cái giá, \(n^2\), và chương~\ref{ch:chi-phi-va-song-song}
đo nó.

\subsection*{Ba chỗ dùng, và chúng không giống nhau}

\index{self-attention!ba chỗ dùng trong kiến trúc}%
Bài dùng attention ở ba chỗ với ba cấu hình khác nhau, và trộn ba chỗ ấy làm
một là cách hiểu sai kiến trúc phổ biến nhất mà tôi gặp. Khác biệt nằm ở query,
key và value lấy từ đâu.

Ba chữ ấy sách này giữ nguyên tiếng Anh, cùng lý do với chữ attention. Bản dịch
sát nhất của chúng là \enquote{truy vấn}, \enquote{khóa} và \enquote{giá trị},
mà \enquote{khóa} và \enquote{giá trị} là hai từ tiếng Việt thường ngày xuất
hiện ở khắp chương này theo nghĩa khác hẳn. Dịch chúng ở đây tốn nhiều hơn được.

\begin{enumerate}
  \item \textbf{Self-attention ở encoder.} Cả ba lấy từ câu nguồn. Mỗi vị trí
    nguồn đọc mọi vị trí nguồn.
  \item \textbf{Self-attention có che ở decoder.} Cả ba lấy từ
    \tn{câu đích}{target sentence}, nhưng
    vị trí \(i\) chỉ được đọc các vị trí từ \(i\) trở về trước.
    Mục~\ref{sec:ch07-che} là chỗ nói vì sao.
  \item \textbf{Attention encoder-decoder.} Query lấy từ decoder, key và value
    lấy từ encoder. \emph{Đây đúng là cơ chế của
    chương~\ref{ch:mot-vector-cho-moi-tu}}, và bài nói thẳng ra thế ở mục 3.2.3:
    nó \enquote{mimics the typical encoder-decoder attention mechanisms in
    sequence-to-sequence models}. Cái mới không phải chỗ này, cái mới là hai
    chỗ kia.
\end{enumerate}

Hình~\ref{fig:ch07-kien-truc} vẽ cả ba.

\begin{figure}[htbp]
  \centering
  % Kiến trúc Transformer, vẽ theo hình 1 của Vaswani và cộng sự 2017.
% Không vẽ đủ N tầng: mỗi ngăn xếp chỉ mở một tầng ra, phần còn lại đánh dấu
% bằng "N x", giống cách hình gốc dùng một khung lặp.
% Trọng tâm của hình không phải liệt kê mọi khối, mà là tách rõ BA chỗ khác
% nhau attention xuất hiện: tự chú ý ở encoder (nhãn 1), tự chú ý có che ở
% decoder (nhãn 2), và attention nối hai phía encoder-decoder (nhãn 3).
% Mono-safe: ba chỗ ấy phân biệt bằng số hiệu, nét đứt/nét chấm và mức xám
% của nền hộp, không dùng màu.
\begin{tikzpicture}[
  >= Stealth,
  hop/.style     = {draw, rectangle, inner sep = 2pt, font = \scriptsize,
                     align = center},
  cong/.style    = {draw, circle, minimum size = 4mm, inner sep = 0pt,
                     font = \scriptsize},
  so/.style      = {draw, circle, fill = white, minimum size = 3.2mm,
                     inner sep = 0pt, font = \tiny},
  diem1/.style   = {fill = black!10},
  diem2/.style   = {densely dashed, line width = 0.8pt},
  diem3/.style   = {densely dotted, line width = 0.8pt, fill = black!22},
  mui/.style     = {->, line width = 0.6pt, black!75},
  phu/.style     = {->, line width = 0.5pt, black!75, densely dashed},
  ngucanh/.style = {->, line width = 1.6pt, black},
  khung/.style   = {black!45, densely dotted, line width = 0.5pt},
  nho/.style     = {font = \scriptsize},
  ghichu/.style  = {font = \footnotesize, black!60},
]

  % ================= ENCODER (cột trái, x = 0) =================
  \draw[mui] (0, -0.7) -- (0, -0.4);
  \node[hop, minimum width = 3.2cm, minimum height = 0.8cm]
    (emb_e) at (0, 0) {Embedding nguồn};

  \draw[mui] (0, 0.4) -- (0, 0.8);
  \node[cong] (pe_e) at (0, 1.0) {\(+\)};
  \node[hop, minimum width = 1.8cm, minimum height = 0.85cm]
    (pos_e) at (-2.0, 1.0) {Mã hóa\\vị trí};
  \draw[mui] (pos_e) -- (pe_e);

  \draw[mui] (0, 1.2) -- (0, 1.7);
  \node[hop, diem1, minimum width = 2.8cm, minimum height = 0.7cm]
    (attn_e) at (0, 2.05) {Self-attention};
  \node[so] at (attn_e.north east) {1};

  \draw[mui] (attn_e) -- (0, 2.8);
  \node[cong] (add1_e) at (0, 3.0) {\(+\)};
  \draw[phu] (attn_e.south) to[out = 160, in = 200] (add1_e.west);
  \node[nho, anchor = west] at (0.35, 3.0) {Cộng và chuẩn hóa};

  \draw[mui] (0, 3.2) -- (0, 3.6);
  \node[hop, minimum width = 2.8cm, minimum height = 0.7cm]
    (ffn_e) at (0, 3.95) {FFN theo vị trí};

  \draw[mui] (ffn_e) -- (0, 4.7);
  \node[cong] (add2_e) at (0, 4.9) {\(+\)};
  \draw[phu] (ffn_e.south) to[out = 160, in = 200] (add2_e.west);

  \draw[mui] (0, 5.1) -- (0, 5.5);

  \begin{scope}[on background layer]
    \draw[khung, rounded corners = 2pt] (-2.3, 1.55) rectangle (2.3, 5.25);
  \end{scope}
  \node[ghichu, anchor = west] at (2.45, 3.4) {N x};

  % ================= DECODER (cột phải, x = 8) =================
  \draw[mui] (8, -0.7) -- (8, -0.4);
  \node[nho, anchor = north] at (8, -0.75) {đầu ra dịch sang phải một bước};
  \node[hop, minimum width = 3.2cm, minimum height = 0.8cm]
    (emb_d) at (8, 0) {Embedding đích};

  \draw[mui] (8, 0.4) -- (8, 0.8);
  \node[cong] (pe_d) at (8, 1.0) {\(+\)};
  \node[hop, minimum width = 1.8cm, minimum height = 0.85cm]
    (pos_d) at (10.0, 1.0) {Mã hóa\\vị trí};
  \draw[mui] (pos_d) -- (pe_d);

  \draw[mui] (8, 1.2) -- (8, 1.7);
  \node[hop, diem2, minimum width = 2.8cm, minimum height = 1.0cm]
    (mattn_d) at (8, 2.2) {Self-attention\\có che};
  \node[so] at (mattn_d.north east) {2};

  \draw[mui] (mattn_d) -- (8, 3.1);
  \node[cong] (add1_d) at (8, 3.3) {\(+\)};
  \draw[phu] (mattn_d.south) to[out = 20, in = 340] (add1_d.east);

  \draw[mui] (8, 3.5) -- (8, 3.9);
  \node[hop, diem3, minimum width = 2.8cm, minimum height = 1.0cm]
    (cattn_d) at (8, 4.4) {Attention\\encoder-decoder};
  \node[so] at (cattn_d.north east) {3};

  \draw[mui] (8, 4.9) -- (8, 5.3);
  \node[cong] (add2_d) at (8, 5.5) {\(+\)};
  \draw[phu] (cattn_d.south) to[out = 20, in = 340] (add2_d.east);

  \draw[mui] (8, 5.7) -- (8, 6.1);
  \node[hop, minimum width = 2.8cm, minimum height = 0.7cm]
    (ffn_d) at (8, 6.45) {FFN theo vị trí};

  \draw[mui] (ffn_d) -- (8, 7.2);
  \node[cong] (add3_d) at (8, 7.4) {\(+\)};
  \draw[phu] (ffn_d.south) to[out = 20, in = 340] (add3_d.east);

  \draw[mui] (8, 7.6) -- (8, 8.0);

  \begin{scope}[on background layer]
    \draw[khung, rounded corners = 2pt] (5.7, 1.55) rectangle (10.3, 7.75);
  \end{scope}
  \node[ghichu, anchor = west] at (10.45, 4.65) {N x};

  % --- Linear và softmax phía trên ngăn xếp decoder
  \draw[mui] (8, 8.0) -- (8, 8.4);
  \node[hop, minimum width = 2.0cm, minimum height = 0.6cm]
    (lin_d) at (8, 8.7) {Linear};
  \draw[mui] (lin_d) -- (8, 9.3);
  \node[hop, minimum width = 2.0cm, minimum height = 0.6cm]
    (sm_d) at (8, 9.6) {softmax};
  \draw[mui] (sm_d) -- (8, 10.3);
  \node[nho, anchor = south] at (8, 10.35) {phân phối xác suất};

  % ================= mũi tên nặng: đầu ra encoder sang attention thứ ba
  % Đây là chỗ duy nhất câu nguồn đi vào decoder, và đúng là cơ chế cả
  % chương 6 đã xây: encoder cấp key và value, decoder chỉ cấp query.
  % Đoạn cuối phải đi ngang vào cạnh tây của ô, không cắt chéo qua nó: một
  % mũi tên 1.6pt đâm vào giữa ô sẽ gạch ngang chính cái nhãn của ô.
  \draw[ngucanh, shorten >= 1pt] (0, 5.5) -- (0, 5.7) -- (5.8, 5.7)
    -- (5.8, 4.4) -- (cattn_d.west);
  \node[nho, anchor = south] at (2.9, 5.78) {đầu ra encoder};

  % ================= chú giải ba điểm attention, đặt dưới cả hai cột để
  % không đè lên hình chính
  \node[nho, anchor = north, align = left, black!70] at (4, -1.9)
    {1: query, key, value đều từ câu nguồn \\
     2: đều từ câu đích, và chỉ được nhìn về sau \\
     3: query từ decoder, key và value từ encoder};

\end{tikzpicture}

  \caption{Kiến trúc, vẽ theo hình 1 của bài, với ba chỗ dùng attention đánh
    dấu riêng. So với hình~\ref{fig:ch06-kien-truc} thì mũi tên đậm vẫn còn và
    vẫn làm đúng việc cũ: nó mang câu nguồn sang cho decoder. Cái mất đi là các
    mũi tên ngang nối bước \(t-1\) với bước \(t\) trong hai nửa. Cái thêm vào là
    ô \tn{mã hóa vị trí}{positional encoding} ở chân mỗi cột, và
    mục~\ref{sec:ch07-vi-tri} nói vì sao bỏ
    recurrence thì bắt buộc phải thêm nó.}
  \label{fig:ch07-kien-truc}
\end{figure}

\subsection*{Cái giá đến ngay lập tức}

\index{self-attention!không phân biệt được thứ tự}%
Bỏ recurrence đi thì mất một thứ mà mạng hồi quy cho không: thứ tự.

Một tổng có trọng số trên một tập hợp không biết tập hợp ấy xếp thế nào. Đảo
thứ tự đầu vào thì đầu ra đảo theo y hệt và không có gì khác đổi. Với
self-attention, một câu và một túi từ là cùng một vật.

Đây là loại phát biểu dễ gật đầu rồi đi tiếp, nên tôi đo. Lấy một tầng
attention, cho một dãy vào, hoán vị dãy ấy, chạy lại, rồi hoán vị ngược kết quả
về và so với lần chạy đầu.

\begin{measured}{sai khác giữa chạy thẳng và chạy trên dãy đã hoán vị}
\begin{minted}{text}
d_model  heads  steps  max|difference|
16       4      6      4.470e-08
32       4      12     4.098e-08
64       8      20     7.451e-08
\end{minted}

Tầng attention một mình, chưa có mã hóa vị trí, trọng số lúc khởi tạo.

Đo bằng \code{experiments/ch07\_position.py} tại tag \code{ch07}.
\end{measured}

Ba con số ấy là nhiễu của số thực dấu chấm động 32 bit, không phải một sai khác
nhỏ. Nghĩa là đẳng thức đúng theo nghĩa đen: tầng này không phân biệt được hai
thứ tự, và không phải nó phân biệt kém.

Bài nói lý do ấy bằng một câu ở đầu mục 3.5, rằng vì mô hình không có
recurrence và không có tích chập nên \enquote{we must inject some information
about the relative or absolute position of the tokens in the sequence}. Bài
không đo. Mục~\ref{sec:ch07-vi-tri} là chỗ trả cái giá này, và nó là mục dài
nhất chương, vì cách trả không hiển nhiên như nó trông.
